\chapter{How to use this book}

\section*{What this book is for}
The interviewer reads your CV for the first time, points at one line, and asks about it. Then they keep digging. This book gets you ready for exactly that. For every line of the Recruitment Agent entry on the Data CV, it gives you the words you need, what to say, what to draw, where the number comes from, and what they will ask next.

\section*{Where it comes from}
The four CV lines come first: where anything else disagrees with them, the CV wins. On top of that: your prep pages (\texttt{timetable/\allowbreak profile/\allowbreak data/\allowbreak recruit/}) and the July--October 2025 code they point to (the LangGraph notebook, the FastAPI backend, the Workspace scripts). The repository's current main branch is a later Django rebuild, used here only for ``what I'd change''.

\section*{The order}
\begin{enumerate}
  \item \textbf{The 30-second version.} Say it when they say ``tell me about this project''.
  \item \textbf{The 3-minute version} (Why, What, How, What now), with the one big drawing.
  \item \textbf{One chapter per CV line}, in CV order. Each chapter has the same parts, in this order:
  \begin{enumerate}
    \item the line, word for word;
    \item \emph{the words you need} to understand it, in plain language and undergraduate maths, before anything uses them;
    \item what to say in 20 seconds;
    \item the drawings (boards);
    \item how it works;
    \item the number and what it really means;
    \item the questions they will ask, with answers;
    \item what not to say.
  \end{enumerate}
  \item \textbf{The honest part}: what the project cannot show yet, and what you would do next.
\end{enumerate}

\section*{The boxes}
\begin{itemize}
  \item Orange: the CV line, word for word.
  \item Blue, numbered: a word explained in plain language. The word is in bold the first time.
  \item Purple, numbered: a worked example. ``Toy'' means made-up numbers that show the idea, not your data.
  \item Dark green: things you say out loud, and the boards to draw on paper.
  \item Green \emph{The number}: where a figure comes from and its honest catch.
  \item Light blue \emph{Q.}: an interview question with the answer to give. A $\bigstar$ marks the ones asked most often.
  \item Red: things that sound good but are wrong or overclaim, and lines to keep fixed.
\end{itemize}
Short on time: read Chapters 1 and 2, then every $\bigstar$ question.

\begin{bugfix}[Have these ready]
\begin{enumerate}
  \item \textbf{Your part.} It was a team project. Know which pieces you built, and say ``we'' for the rest.
  \item \textbf{The 80\%.} No timing log is saved. Your fixed definition: recruiter hands-on minutes per role, from request to invites sent, manual against agent. Have your own breakdown ready (Chapter 4 has an illustrative one).
  \item \textbf{The 500+.} Applications ranked, added up over every role you ran it on. Know roughly how many roles.
\end{enumerate}
\end{bugfix}
